
\subsection{Desempe\~no Computacional}

Para evaluar el rendimiento del modelo y la eficiencia del motor \textbf{lp\_solve}, se registraron las dimensiones matem\'aticas estructurales (cantidad de variables booleanas y restricciones) junto con el tiempo de CPU en segundos para el conjunto de 15 instancias de prueba. Todas las ejecuciones procesadas alcanzaron exitosamente el estado de \textbf{\'Optimo} global, demostrando la factibilidad y completitud del modelo. Los resultados emp\'iricos se resumen en el Cuadro \ref{tab:desempeno}.

\begin{table}[htbp]
\caption{M\'etricas de Desempe\~no Computacional por Instancia}
\label{tab:desempeno}
\centering
\begin{tabular}{|l|c|c|c|r|c|}
\hline
\textbf{Instancia} & \textbf{Tipo} & \textbf{Var.} & \textbf{Restr.} & \textbf{Tiempo (s)} & \textbf{Costo Total (\$)} \\
\hline
Peque\~na\_1  & Peque\~na & 39  & 5  & 0.0050 & 14,116,009.97 \\
Peque\~na\_2  & Peque\~na & 22  & 5  & 0.0056 & 17,986,055.80 \\
Peque\~na\_3  & Peque\~na & 36  & 5  & 0.0046 & 27,157,527.11 \\
Peque\~na\_4  & Peque\~na & 22  & 4  & 0.0046 & 20,073,387.31 \\
Peque\~na\_5  & Peque\~na & 16  & 3  & 0.0045 & 15,760,216.80 \\
\hline
Mediana\_6   & Mediana & 109 & 8  & 0.0099 & 24,833,358.63 \\
Mediana\_7   & Mediana & 97  & 6  & 0.0060 & 11,639,670.16 \\
Mediana\_8   & Mediana & 70  & 6  & 0.0071 & 20,644,785.16 \\
Mediana\_9   & Mediana & 117 & 8  & 0.0112 & 20,588,708.87 \\
Mediana\_10  & Mediana & 116 & 7  & 0.0058 & 28,011,806.04 \\
\hline
Grande\_11   & Grande  & 427 & 12 & 0.0444 & 31,060,548.51 \\
Grande\_12   & Grande  & 502 & 15 & 0.0995 & 33,180,921.26 \\
Grande\_13   & Grande  & 471 & 14 & 0.1681 & 12,947,800.61 \\
Grande\_14   & Grande  & 478 & 14 & 0.0687 & 43,557,487.56 \\
Grande\_15   & Grande  & 513 & 14 & 0.0876 & 23,577,441.11 \\
\hline
\end{tabular}
\end{table}

Dado que el motor confirm\'o la optimalidad en el 100\% de las pruebas, podemos asegurar que la brecha de optimalidad (MIP gap) es nula (0\%). Adem\'as, la respuesta fue tan r\'apida que el algoritmo no tuvo necesidad de explorar ramificaciones profundas; resolvi\'o los escenarios pr\'acticamente de entrada, sin sobrecargar la capacidad de procesamiento.\\

Cuando observamos las instancias Peque\~nas y Medianas, vemos un entorno sumamente ligero (hasta 117 variables y 8 restricciones). Aqu\'i, la m\'aquina realiz\'o el trabajo de forma inmediata, demorando apenas entre $0.0045$ y $0.0112$ segundos.\\

Al escalar a las instancias Grandes, el desaf\'io l\'ogicamente crece, llegando a manejar hasta 513 variables en los casos m\'as complejos. Sin embargo, el comportamiento m\'as revelador para el equipo lo encontramos en la instancia \texttt{Grande\_13}. Aunque es estructuralmente m\'as peque\~na que la \texttt{Grande\_12} (471 variables y 14 restricciones, frente a 502 y 15), le tom\'o casi un 70\% m\'as de tiempo llegar a la soluci\'on ($0.1681$ segundos).\\

¿Por qu\'e ocurre esto? Nos demuestra que el tama\~no en bruto no es lo \'unico que dicta el rendimiento. Cuando las exigencias de aislaci\'on t\'ermica ($U^{\text{max}}$) est\'an al l\'imite, o cuando los costos de los materiales son demasiado similares entre s\'i, el motor se ve obligado a analizar detalladamente muchas m\'as posibilidades antes de poder garantizarle al usuario que encontr\'o la opci\'on m\'as econ\'omica.\\

Aun con esta variabilidad, ning\'un escenario super\'o los dos d\'ecimos de segundo de espera. Para el equipo, esto confirma emp\'iricamente que nuestra decisi\'on de dise\~no m\'as importante —precalcular los componentes fijos por fuera y entregarle al motor \'unicamente las decisiones libres— fue un acierto. Logramos crear un sistema inteligente y ágil, que no desperdicia recursos y que est\'a listo para escalar en proyectos reales.
